\section{Auditoria, Rastreabilidade e Versionamento}
\label{sec:auditoria-rastreabilidade-versionamento}

Esta seção reúne os conceitos que sustentam a principal contribuição do módulo: a capacidade
de registrar, consultar e desfazer alterações cadastrais com garantias de rastreabilidade.

\subsection{Auditoria e Rastreabilidade}
\label{subsec:auditoria-rastreabilidade}

Auditoria, no contexto de sistemas de informação, é o registro sistemático de quem realizou
determinada operação, quando, sobre qual dado e com qual resultado. A rastreabilidade é a
propriedade que permite reconstituir o histórico de uma informação a partir desses registros.
Fowler \cite{fowler2004} descreve o padrão \textit{Audit Log} como um registro de alterações
feitas nos dados, com o objetivo de permitir verificar o que ocorreu e, quando necessário,
compreender a origem de um estado atual. Em sistemas que envolvem decisões de classificação e
seleção de pessoas, como o Banco de Especialistas, esse registro contribui para a transparência
e para a responsabilização dos agentes envolvidos.

O fluxo anterior ao módulo, baseado em e-mail e execução manual de SQL, ilustra a ausência
dessas propriedades: não havia registro formal do solicitante, da justificativa, do responsável
pela decisão ou do estado anterior do cadastro. Ao formalizar a solicitação como uma entidade
persistente, com status, autoria e datas, o sistema passa a produzir esse registro como efeito
natural do próprio fluxo de trabalho.

% #TO-DO: Ilustração que represente o fluxo anterior baseado em e-mails em constraste com o fluxo
% formalizado pelo módulo.

\subsection{Fluxo de Aprovação com Máquina de Estados}
\label{subsec:fluxo-aprovacao-maquina-estados}

Uma solicitação percorre um ciclo de vida bem definido, que pode ser modelado como uma máquina
de estados finitos: um conjunto de estados (por exemplo, pendente, aprovada e rejeitada) e de
transições permitidas entre eles \cite{sommerville2011}. Explicitar as transições válidas evita
estados incoerentes: uma solicitação rejeitada não pode ser aprovada sem nova análise, e somente
uma solicitação aprovada pode ser revertida. A validação dessas regras no servidor garante que o
fluxo seja respeitado independentemente da interface utilizada.

% #TO-DO: diagrama UML de máquina de estados, mostrando estados reais do sistema.

\subsection{Snapshot e Reversão}
\label{subsec:snapshot-reversao}

Para possibilitar a reversão de uma alteração, é necessário preservar o estado anterior dos
dados. O conceito de \textit{snapshot} (instantâneo) consiste em capturar e armazenar o estado
de um objeto em determinado momento. Ele se aproxima do padrão \textit{Memento}, descrito por
Gamma et al. \cite{gamma1994}, cujo propósito é capturar o estado interno de um objeto, sem
violar seu encapsulamento, para que ele possa ser restaurado posteriormente.

No módulo, a entidade \textit{Snapshot} registra o estado relevante do especialista (como
categoria e pontuação) antes da aplicação de uma alteração aprovada. A reversão consiste em
restaurar esses valores, mantendo a consistência entre categoria e pontuação. Cabe distinguir
essa abordagem da simples exclusão do registro de histórico: a reversão é, ela mesma, uma
operação registrada, de modo que o histórico permanece completo e a operação inversa pode ser
auditada. Esse comportamento é análogo à ideia de transações compensatórias, em que o efeito
de uma operação já efetivada é desfeito por uma nova operação, e não pela remoção do
passado.

Existem abordagens mais gerais para o versionamento de dados, como o \textit{Event Sourcing}, em
que o estado é derivado de uma sequência imutável de eventos \cite{fowler2005, kleppmann2017}. Embora o
Event Sourcing seja uma alternativa robusta, o módulo adota uma solução mais simples e proporcional ao
escopo, combinando \textit{snapshots} do estado anterior com um registro de eventos, e o conceito
de \textit{Event Sourcing} é aqui mencionado como referência conceitual para a ideia de histórico imutável.

% TO-DO: Diagrama temporal que explique o antes, alteração e a restauração.

\subsection{Evento de Domínio e Domain-Driven Design}
\label{subsec:evento-dominio-ddd}
O Domain-Driven Design (DDD) é uma abordagem de design de software que coloca o domínio do
negócio no centro do desenvolvimento, utilizando uma linguagem ubíqua para alinhar especialistas
de negócio e desenvolvedores \cite{evans2003}. Dentro do DDD, o conceito de Evento de Domínio
(Domain Event) representa algo significativo que aconteceu no domínio do problema. Um evento de
domínio é, por definição, um registro imutável de um fato passado (por exemplo, "Solicitação
Aprovada" ou "Categoria Alterada"), que captura a intenção e o contexto da operação
\cite{evans2003}.

No contexto do Banco de Especialistas, o módulo de revisão foi modelado sob essa ótica. A
entidade SolicitationReview não atua apenas como uma tabela de auditoria, mas como o registro do
Evento de Domínio que materializa a decisão tomada pelo avaliador. Ao registrar o evento (a
revisão) de forma imutável e com autoria definida, o sistema garante que as transições de estado
do especialista (como a mudança de categoria) sejam disparadas a partir de um fato consumado e
rastreável, promovendo o desacoplamento entre a intenção (solicitação) e a sua efetivação
(evento de aprovação).

\subsection{Timeline e Registro de Eventos}
\label{subsec:timeline-registro-eventos}

A \textit{timeline} é a apresentação ordenada cronologicamente dos eventos associados a uma
entidade. Cada evento (criação, aprovação, rejeição, reversão) é registrado com seu tipo,
autor, data e informações complementares, e a interface os exibe como uma linha do tempo.
Trata-se de uma forma de tornar a auditoria acessível ao usuário final, e não apenas ao
administrador do banco de dados.

% TO-DO: Imagem da tela da timeline mostrando diversos eventos de uma solicitação, 
% incluindo a reversão de uma alteração.
